% 共用导言：三册书共用的格式定义
% 编译方式：xelatex（需 TeX Live + ctex + tcolorbox）

% 注意：ctex 的选项（fontset=none, UTF8, scheme=chinese）在 \documentclass 处给出
\usepackage{ctex}
\setCJKmainfont{Noto Serif CJK SC}
\setCJKsansfont{Noto Sans CJK SC}
\setCJKmonofont{Noto Sans Mono CJK SC}
% 西文使用默认 Latin Modern，与 Noto CJK 搭配

\usepackage[a4paper,top=2.6cm,bottom=2.6cm,left=2.4cm,right=2.4cm]{geometry}
\usepackage{amsmath,amssymb,bm}
\usepackage{graphicx}
\usepackage{xcolor}
\usepackage{tikz}
\usepackage{booktabs,longtable,array,multirow}
\usepackage{enumitem}
\usepackage{zhnumber}
\usepackage[most]{tcolorbox}
\usepackage[hidelinks]{hyperref}

% ---------- 颜色语法（全书统一的人为编码，颜色不代表真实外观） ----------
\definecolor{mainblue}{RGB}{31,78,121}     % 主线
\definecolor{bridgegreen}{RGB}{46,117,80}  % 原理桥梁
\definecolor{challengepurple}{RGB}{96,60,130} % 深入挑战
\definecolor{laborange}{RGB}{176,96,32}    % 实验
\definecolor{warnred}{RGB}{158,42,43}      % 误解警示
\definecolor{evidenceteal}{RGB}{30,110,110}% 证据故事

% ---------- 化学式快捷写法：\chem{H_2O}、\chem{CO3^{2-}} ----------
\newcommand{\chem}[1]{\ensuremath{\mathrm{#1}}}

% ---------- 三层阅读标记 ----------
% 故事主线是正文默认层，无需标记；以下两层用色块标出
\newtcolorbox{qiaoliang}{% 原理桥梁：定量规律
  enhanced,breakable,colback=bridgegreen!6,colframe=bridgegreen,
  boxrule=0.6pt,arc=1.5mm,left=2mm,right=2mm,
  title={\sffamily 【原理桥梁】用一点数学看得更深},
  fonttitle=\bfseries,coltitle=white,colbacktitle=bridgegreen}

\newtcolorbox{tiaozhan}{% 深入挑战：更高层的视野
  enhanced,breakable,colback=challengepurple!6,colframe=challengepurple,
  boxrule=0.6pt,arc=1.5mm,left=2mm,right=2mm,
  title={\sffamily 【深入挑战】想再往上走一层吗},
  fonttitle=\bfseries,coltitle=white,colbacktitle=challengepurple}

% ---------- 固定栏目 ----------
\newtcolorbox{kaichang}{% 开场物品与悬念
  enhanced,breakable,colback=mainblue!5,colframe=mainblue,
  boxrule=0.6pt,arc=1.5mm,left=2mm,right=2mm,
  title={\sffamily 开场：先拿起一样东西},
  fonttitle=\bfseries,coltitle=white,colbacktitle=mainblue}

\newtcolorbox{shiyan}{% 安全的家庭/课堂观察实验
  enhanced,breakable,colback=laborange!6,colframe=laborange,
  boxrule=0.6pt,arc=1.5mm,left=2mm,right=2mm,
  title={\sffamily 动手试一试},
  fonttitle=\bfseries,coltitle=white,colbacktitle=laborange}

\newtcolorbox{wuqu}{% 常见误解与反例
  enhanced,breakable,colback=warnred!5,colframe=warnred,
  boxrule=0.6pt,arc=1.5mm,left=2mm,right=2mm,
  title={\sffamily 小心：这里有个常见误解},
  fonttitle=\bfseries,coltitle=white,colbacktitle=warnred}

\newtcolorbox{zhengju}{% 科学家怎样知道
  enhanced,breakable,colback=evidenceteal!5,colframe=evidenceteal,
  boxrule=0.6pt,arc=1.5mm,left=2mm,right=2mm,
  title={\sffamily 科学家怎样知道},
  fonttitle=\bfseries,coltitle=white,colbacktitle=evidenceteal}

% 想一想（思考题）
\newlist{sikaoti}{enumerate}{1}
\setlist[sikaoti,1]{label=\textbf{想一想 \arabic*.},leftmargin=3.2em,itemsep=2pt}

% ---------- 章节样式 ----------
\ctexset{
  chapter = {
    name = {第,章},
    number = \chinese{chapter},
    format = \Huge\bfseries\sffamily,
  },
  section = {format = \Large\bfseries\sffamily},
  subsection = {format = \large\bfseries\sffamily},
}

% 篇（part）命名
\ctexset{
  part = {
    name = {第,篇},
    number = \chinese{part},
    format = \Huge\bfseries\sffamily\centering,
  }
}

% 页眉页脚
\usepackage{fancyhdr}
\pagestyle{fancy}
\fancyhf{}
\fancyhead[LE]{\small\sffamily 烧掉化学与生物课本}
\fancyhead[RO]{\small\sffamily \leftmark}
\fancyfoot[C]{\small\thepage}
\renewcommand{\headrulewidth}{0.4pt}

% 行距与段落
\linespread{1.32}
\setlength{\parskip}{2pt}

% 常用物理量与单位快捷命令
\newcommand{\ud}[2]{\ensuremath{#1\ \mathrm{#2}}} % \ud{6.02\times10^{23}}{mol^{-1}}
